\section{尚哥指标清单的证据可用性矩阵}
\label{app:metric-matrix}

本附录逐项对应尚哥提供的模型、数据和指标提纲。全文涉及的五层分母严格分开：
\begin{enumerate}
  \item \textbf{单体：}Ser 来源修正留出集，151条、1505个真实位点，正文阈值严格 $p>0.6$；
  \item \textbf{六复合物主分析：}六个非 3AV 靶点、$T=0.5$、544条原始序列和476条甲基化保留候选；
  \item \textbf{天然17靶点：}预先选定的全部短肽链，共251个天然位点；
  \item \textbf{17靶点全温度生成池：}五个温度共4015条原始记录、3443条温度内唯一序列；
  \item \textbf{best85 探索集：}17个靶点、五个温度、按天然序列恢复率选择的85个 oracle 代表。
\end{enumerate}

状态含义如下：\textbf{可报告}表示定义、分母和来源均已冻结；\textbf{探索性}表示数值存在，但受 oracle 选择、预测模型或非实验代理限制；\textbf{不可定义}表示真值类别缺失，数学上不能计算；\textbf{未计算}表示没有合格的成对数据或冻结定义。任何“不可定义/未计算”均不等于数值0。

\subsection{模型、数据与序列指标}

\begin{landscape}
\begingroup
\small
\setlength{\tabcolsep}{3pt}
\renewcommand{\arraystretch}{1.22}
\begin{longtable}{L{27mm}L{46mm}L{50mm}L{48mm}L{43mm}}
\caption{模型、数据与序列指标的队列级可用性。}\label{tab:metric-matrix-sequence}\\
\toprule
清单项目 & 单体 & 六复合物主分析 & \makecell{17靶点探索\\（天然/生成/best85）} & 状态与原因 \\
\midrule
\endfirsthead
\multicolumn{5}{l}{\small\itshape 表\thetable（续）}\\
\toprule
清单项目 & 单体 & 六复合物主分析 & \makecell{17靶点探索\\（天然/生成/best85）} & 状态与原因 \\
\midrule
\endhead
\midrule
\multicolumn{5}{r}{\small 接下页}\\
\endfoot
\bottomrule
\endlastfoot

基于甲基化单体的微调
& 冻结检查点可评价；历史审计记录 train/test 为600/151，但无原始训练全过程日志、随机种子曲线和相对于初始化模型的配对增益
& 使用同一冻结检查点生成/评分，不构成新的训练证据
& 天然集、全温度池和 best85 均使用同一冻结检查点；oracle 选择不能反推微调增益
& \textbf{部分可报告}：可写模型和检查点，不可编造“微调提高了多少” \\

甲基化分类器
& known-sequence 与 end-to-end 均有261阳性/1244阴性真值
& 75个天然位点全为阴性；476条生成序列无实验甲基真值
& 天然251位点全为阴性；全温度池和85条代表只有模型小写标记，无实验甲基真值
& 单体\textbf{可报告}；复合物生成序列只能报告标记/保留率 \\

Baker 33 个真结构单体
& 未找到；751/751 个原始 PDB 均标记为 Rosetta 理论模型，没有可核验的 Baker 实验结构子集
& 不适用
& 不适用
& \textbf{与冻结数据不符}：不得把规划数字写成已完成数据 \\

Rosetta 400+300 个单体
& 实际为406个 Rosetta-2023 与345个 Rosetta-2025 理论模型，合计751；训练/测试为600/151
& 不适用
& 不适用
& \textbf{可报告实际构成}：规划的400+300数字不采用 \\

17 个真结构复合物
& 不适用
& 从17个靶点中预设选取六个非3AV靶点
& 天然17靶点用于序列评价；best85为85/85结构匹配
& \textbf{可报告}，但六靶点与 best85 的队列和 RMSD 定义不同 \\

RAA（天然化氨基酸恢复）
& $242/1505=16.08\%$
& 天然六复合物七条选定肽链为 $29/75=38.67\%$；不得与476条生成候选的回收字段混为同一测试
& 天然17靶点为$53/251=21.12\%$；全温度池按温度的平均RAA为1.49--4.42\%；best85的 recovery 是 oracle 选择标准
& 单体、天然六复合物和天然17靶点\textbf{可报告}；生成池为描述性，best85有\textbf{选择偏倚} \\

AUC
& known-sequence 0.9003；end-to-end 0.9082
& 天然75位点阳性数为0，AUC不可定义；生成476条无实验真值
& 天然251位点无阳性，AUC不可定义；生成池与best85无实验甲基真值
& 仅单体\textbf{可报告} \\

Accuracy（严格 $p>0.6$）
& known $1297/1505=86.18\%$；end-to-end $1341/1505=89.10\%$
& 天然阴性集 known $71/75=94.67\%$；end-to-end $73/75=97.33\%$，实质为 $1-\mathrm{FPR}$；生成476条不可计算
& 天然17靶点 known $232/251=92.43\%$、end-to-end $242/251=96.41\%$，实质为$1-\mathrm{FPR}$；生成池与best85不能计算
& 单体\textbf{可报告}；两个天然复合物阴性集仅作误报审计 \\

Precision / Recall / F1
& known 64.02\%/46.36\%/53.78\%；end-to-end 88.80\%/42.53\%/57.51\%
& 天然位点无阳性，三项不可定义；生成候选无真值
& 天然251位点无阳性，三项不可定义；生成池与best85无实验真值
& 仅单体\textbf{可报告} \\

RMC
& 二分类甲基召回 $111/261=42.53\%$；端到端真实甲基残基身份完全恢复 $20/261=7.66\%$；全部位点联合 token 恢复 $224/1505=14.88\%$
& 同理；因无阳性，不能把RMC写成0
& 天然251位点无阳性，不能把RMC写成0；生成池与best85无真值
& \textbf{术语歧义}：必须同时写明身份要求与分母 \\

FPR / 预测阳性率
& known FPR 5.47\%、预测阳性率12.56\%；end-to-end 1.13\%、8.31\%
& 天然阴性集 known $4/75=5.33\%$；end-to-end $2/75=2.67\%$
& 天然17靶点 known FPR $19/251=7.57\%$；end-to-end预测阳性率$9/251=3.59\%$；生成池可统计小写率但不能称FPR
& 单体及两个天然阴性集\textbf{可报告} \\

联合扩展 token 恢复
& 基础类别与甲基状态同时正确：$224/1505=14.88\%$
& 未形成可复核的75位点联合交叉表；生成候选无扩展 token 真值
& 天然17靶点未形成联合交叉表；生成池与best85无扩展 token 真值
& 仅单体\textbf{可报告}；不能用 RAA 与 Accuracy 相乘替代 \\

甲基化保留率/有效生成概率
& 不适用；这是生成队列指标
& $476/544=87.50\%$ 含至少一个严格 $p>0.6$ 的已保存小写标记；联合 $<3$/$<5$ 原始产率为 $16/544=2.94\%$、$101/544=18.57\%$
& 全温度池有4015条记录，但未定义与主队列相同的端到端结构成功门；best85固定为85个oracle代表
& 六复合物\textbf{主结果可报告} \\

旧标签/旧 padding 阈值截图
& 不进入正文；旧标签含62个Ser来源错误，旧截图分母随 batch size 改变
& 不适用
& 不适用
& \textbf{仅历史审计}：不可与本表正文指标并列 \\

\end{longtable}
\endgroup
\end{landscape}

\subsection{结构、能量、多样性与通透性指标}

\begin{landscape}
\begingroup
\footnotesize
\setlength{\tabcolsep}{3pt}
\renewcommand{\arraystretch}{1.20}
\begin{longtable}{L{29mm}L{43mm}L{51mm}L{52mm}L{40mm}}
\caption{结构与下游指标的队列级可用性。}\label{tab:metric-matrix-structure}\\
\toprule
清单项目 & 单体 & 六复合物主分析 & \makecell{17靶点探索\\（天然/生成/best85）} & 状态与原因 \\
\midrule
\endfirsthead
\multicolumn{5}{l}{\small\itshape 表\thetable（续）}\\
\toprule
清单项目 & 单体 & 六复合物主分析 & \makecell{17靶点探索\\（天然/生成/best85）} & 状态与原因 \\
\midrule
\endhead
\midrule
\multicolumn{5}{r}{\small 接下页}\\
\endfoot
\bottomrule
\endlastfoot

RMSD\_CA
& 无端到端单体设计结构--参考结构配对
& 476/476有一次全复合物对齐后的全局 CA 与末链 best-forward 环肽 CA RMSD；总体中位数[IQR]分别为1.022[0.692--1.534]和6.738[5.261--8.129]~\AA
& 受体 CA 拟合后肽 CA mean/median=29.93/40.06~\AA；肽自对齐=2.65/2.53~\AA
& 六复合物与 best85 均\textbf{可报告}，但定义不同，禁止合并 \\

RMSD\_Backbone
& 未计算
& 476条主队列未计算完整 backbone RMSD；只有选择后的旧代表字段，不能替代全队列
& N/CA/C：受体拟合后 mean/median=29.83/39.96~\AA；肽自对齐=2.34/2.11~\AA，85/85
& best85\textbf{可报告}；六主队列\textbf{未计算} \\

RMSD\_All
& 未计算
& 未计算
& 未计算
& \textbf{缺失}：需要非标准/甲基残基的可靠全原子映射 \\

\makecell[l]{RMSD\_CA\\Methylation}
& 未计算
& 六个按肽 RMSD 选择的代表有旧位点字段，但不是476条主队列的无偏汇总
& 159个甲基标记位点与626个非甲基位点：受体拟合 pooled CA=20.43/34.18~\AA；肽自拟合=3.28/2.93~\AA
& best85\textbf{描述性}；分组非随机，不能推断甲基化因果效应 \\

\makecell[l]{RMSD\_Backbone\\Methylation}
& 未计算
& 仅选择后的旧代表字段；主队列未汇总
& N/CA/C 位点结果按温度可用，85/85；无显式甲基原子
& best85\textbf{描述性} \\

\makecell[l]{RMSD\_All\\Methylation}
& 未计算
& 未计算
& 未计算
& \textbf{缺失}：没有显式 N-甲基全原子对应关系 \\

pLDDT
& 无单体设计结构结果
& 六个选择后代表仅有 \path{highfold_ca_bfactor_mean}，不能改名为 COMMENT global pLDDT
& COMMENT global pLDDT 37/85（mean 71.07）；CA B-factor代理85/85（mean87.55），二者分开报告
& best85\textbf{部分可用}；缺失不得填0 \\

ipLDDT
& 未定义/未计算
& 未定义/未计算
& 没有名为 ipLDDT 的冻结字段；最近的字段是 peptide-chain pLDDT，37/85（mean46.53）
& \textbf{术语未定义}：若指界面/肽链pLDDT，须按真实字段改名 \\

pTM
& 未计算
& 六个选择后代表缺失
& peptide-chain pTM 37/85（mean0.146）
& best85\textbf{仅部分可用} \\

ipTM
& 未计算
& 六个选择后代表缺失
& peptide--receptor ipTM 37/85（mean0.461）
& best85\textbf{仅部分可用} \\

pAE / PAE
& 未计算
& 六个选择后代表缺失
& peptide-chain PAE mean3.64、inter-PAE mean18.05、iPAE mean12.14，均37/85
& best85\textbf{仅部分可用}；不能由PDB坐标补造 \\

BSR（binding-site ratio）
& 未定义/未计算
& 未定义/未计算
& 未定义/未计算
& \textbf{缺失}：须先冻结接触原子、距离阈值以及 precision/recall 方向 \\

Success：肽能量低于天然
& 无设计/天然同流程能量对
& 六代表只有设计 fixed-pose 分数，无同流程天然肽基线
& 只有设计 fixed-pose 分数，无同流程天然肽基线
& \textbf{不可计算}；界面分数 $<0$ 不能替代“低于天然” \\

Stability：复合物能量低于天然
& 无设计/天然同流程能量对
& 无同流程天然复合物能量
& 无同流程天然复合物能量
& \textbf{不可计算} \\

Designability：RMSD $<2$~\AA
& 无结构对
& 按 Task 1 联合定义为 $2/476=0.42\%$；以全部544条为分母为 $2/544=0.37\%$
& 受体框架肽 CA RMSD $<2$ 为 $0/85$
& \textbf{可由现有RMSD计算}；须连同定义和分母报告 \\

预设 RMSD $<3$/$<5$~\AA
& 不适用
& 保留后联合通过 $16/476=3.36\%$、$101/476=21.22\%$；原始产率2.94\%、18.57\%
& 旧受体框架肽 CA $<5$ 为 $3/85=3.53\%$，$<3$ 为 $0/85$
& \textbf{可报告}，但两队列定义不同 \\

Diversity：$1-\mathrm{TM}$
& 未计算
& 每靶点只取一个选择后代表，靶内多样性不可估计
& 17靶点内170/170对；mean/median $1-\mathrm{TM}=0.660/0.691$
& best85\textbf{探索性}：肽自对齐内部形状且受oracle选择影响 \\

膜通透性
& 无单体结果
& 六个按肽 RMSD 选择的代表6/6有预测值；属于选择条件化展示，不代表476条总体
& best85为85/85有预测值，其中82条严格温度匹配、3条相同序列跨温度回填；总体 median $1.84\times10^{-8}$，无实验标签
& \textbf{探索性预测}：不能声称实验透膜性提高 \\

Rosetta 能量/分数
& 无单体结果
& 六个选择后代表有自然化 fixed-pose 界面分数；仅作条件化示例
& 85/85；complex score/residue median1.7995 REU，cross-interface median110.07 REU，$<0$ 为6/85
& \textbf{探索性}：非自由能、非亲和力、非相对天然稳定性 \\

单体结构结果总体状态
& 与本文阈值0.60一致的设计结构和成对结果表缺失；旧阈值任务清单不能用于本文
& 不适用
& 复合物结果不得移植到单体
& \textbf{明确缺失}：单体结构指标均不填数 \\

\end{longtable}
\endgroup
\end{landscape}

\subsection{矩阵所引用的冻结数据资产}

上述数值来自随稿提供的单体逐位主表与标准化阈值曲线、476条候选总表、六复合物主表以及 best85 的结构、甲基位点、多样性、通透性和 fixed-pose 能量 CSV。机器可读文件统一位于 \path{data/}，文件索引、关键哈希与上游逻辑来源见 \path{SOURCE_MANIFEST.json}。旧终端截图仅保留在补充材料中，用于说明 padding 与历史标签错误；其数值不参与正文、摘要或本矩阵的主结果。
